\begin{corollary}[7.6.6]
\label{IV.2.s7.6.6-fr}
Sean $A$ un anillo semilocal noetheriano que es un dominio de integridad de
dimensión 1 y $K$ su cuerpo de fracciones. Para que $A$ sea un anillo japonés,
es necesario y suficiente que el completado $\widehat A$
\oldpage[IV-2]{212}
de $A$ sea reducido y que, si $\mathfrak q_j$ \textup{(}$1\leq j\leq n$\textup{)} son los
ideales primos minimales de $\widehat A$, los cuerpos de fracciones de los
dominios de integridad $\widehat A/\mathfrak q_j$ sean extensiones separables
de $K$.
\end{corollary}

\begin{proof}
En efecto, los cocientes de $A$ que son dominios de integridad son el propio
$A$ y cuerpos; por tanto, la condición $b)$ de (7.6.4) equivale a que $A$ sea
un anillo japonés, y la condición $a')$ equivale a las condiciones del
enunciado relativas a $\widehat A$. De aquí se sigue la conclusión.
\end{proof}
